\documentclass{templates/forecastp23}

\newcommand{\authorinfo}{THE FORUM CENTER - NGUYỄN HOÀNG HUY}
\newcommand{\documenttitle}{TÀI LIỆU IELTS SPEAKING QUÝ 3}
\newcommand{\subtitle}{IELTS SPEAKING PART 2 + 3}

\begin{document}

\thispagestyle{firstpage}
\makeforecastheader{\authorinfo}{\documenttitle}{\subtitle}
\pagestyle{otherpages}

\begin{center}{\color{topicblue}\sffamily\bfseries Topic 1: Mobile Phones}\end{center}
\vspace{10pt}
\forecastcuecard{Describe an occasion when you were not allowed to use your mobile phone}{%
  \item \textbf{When it was}
  \item \textbf{Where it was}
  \item \textbf{Why you were not allowed to use your mobile phone}
  \item And how you felt about it
}

\forecastsample[1]{You know how glued we are to our phones these days? Well, I had a pretty intense reality check regarding that dependency last December during my final exams. Specifically, it was the morning of my Advanced Mathematics paper, which is notorious for being the hardest exam of the year.\par The university administration has always been strict, but this time they took it to a whole new level. Walking into the main creative hall, which usually feels quite welcoming, felt more like entering a high-security facility. There were signs everywhere explicitly stating "Zero Tolerance," and we had to surrender all electronics into these sealed bags at the entrance. They even had signal jammers active, just to make sure nobody could try anything clever with a smartwatch or earpiece.\par At first, the silence in my pocket was deafening. I honestly felt this phantom vibration every few minutes, my hand instinctively twitching towards my hip whenever I hit a mental roadblock with a calculus problem. It was almost embarrassing to realize how much I rely on that little black mirror as a digital pacifier to soothe my anxiety.\par But here’s the strange thing—about twenty minutes in, that anxiety just vanished. Without the subconscious part of my brain constantly waiting for a notification or wondering what was happening on Instagram, I entered this state of deep flow that I hadn't experienced in months. My mind felt clearer, sharper. I wasn't just solving equations; I was actually enjoying the process. Walking out of that hall three hours later, the first thing I did wasn't to check my messages like I usually would. I actually just stood there for a moment, enjoying the winter air, realizing that disconnecting might be exactly what I need to survive the rest of this degree.}
\clearpage
\begin{forecastvocabtable}
\forecastvocabrow{Reality check}{n}{/riˈæləti tʃek/}{an occasion that makes you consider the facts exactly as they are}{sự thức tỉnh trước thực tế}
\forecastvocabrow{Notorious}{adj}{/nəʊˈtɔːriəs/}{famous for something bad}{khét tiếng}
\forecastvocabrow{Zero Tolerance}{n}{/ˌzɪərəʊ ˈtɒlərəns/}{strict enforcement of regulations}{không khoan nhượng}
\forecastvocabrow{Deafening}{adj}{/ˈdefənɪŋ/}{(metaphorically) very loud or noticeable}{(nghĩa bóng) đinh tai nhức óc}
\forecastvocabrow{Phantom vibration}{n phrase}{/ˈfæntəm vaɪˈbreɪʃn/}{feeling a phone vibrate when it didn't}{rung ảo}
\forecastvocabrow{Digital pacifier}{n phrase}{/ˈdɪdʒɪtl ˈpæsɪfaɪə(r)/}{something digital used to calm oneself}{"núm vú giả" kỹ thuật số}
\forecastvocabrow{Deep flow}{n phrase}{/diːp fləʊ/}{a mental state of complete focus}{trạng thái tập trung sâu}
\end{forecastvocabtable}

\clearpage
\begin{center}{\color{topicblue}\sffamily\bfseries Mobile Phones}\end{center}
\vspace{10pt}
\forecastqa{How do young and old people use mobile phones differently?}{There are quite stark differences, I'd say. Young people like myself tend to use phones as all-in-one devices – we're constantly on social media, streaming content, and using various apps for everything from food delivery to banking. We're also more comfortable with new features and updates. Older generations, from what I've observed with my parents and relatives, typically use phones more conservatively – mainly for calls, texts, and perhaps basic apps like Zalo or Facebook. They're often less inclined to explore new functions and tend to stick with what they know works. There's also a difference in dependency – younger people seem almost addicted to their devices, while older folks can easily put their phones aside for hours.}
\forecastqa{What positive and negative impact do mobile phones have on friendship?}{On the positive side, mobile phones have made it incredibly easy to stay connected with friends, even when we're physically apart. I can maintain friendships with people from my hometown or even international friends I've met online. Group chats keep everyone in the loop, and we can share moments instantly through photos and videos. However, there's definitely a downside. I've noticed that when friends get together, we sometimes end up scrolling through our phones rather than having meaningful face-to-face conversations. There's also this phenomenon where people present idealized versions of their lives on social media, which can create superficial comparisons and even jealousy among friends. The quality of interaction has arguably diminished despite the increased quantity of contact.}
\forecastqa{Is it a waste of time to take pictures with mobile phones?}{I wouldn't say it's a waste of time, but it depends on the approach. Taking photos to capture meaningful moments and memories is perfectly reasonable – I love looking back at pictures from trips or gatherings with friends. However, it becomes problematic when people are so focused on getting the perfect shot for Instagram that they're not actually experiencing the moment itself. I've been guilty of this myself, spending ages trying to photograph my food instead of enjoying it while it's hot. The key is finding a balance – document your life, but don't let the documentation become more important than the actual living.}
\forecastqa{Do you think it is necessary to have laws on the use of mobile phones?}{Absolutely, in certain contexts. Laws banning phone use while driving are essential – I've seen too many near-accidents caused by distracted drivers. Similarly, regulations about phone use in schools or during exams make perfect sense to maintain focus and prevent cheating. However, I think overly restrictive laws that try to control phone use in all aspects of life would be excessive and difficult to enforce. It's more about targeting specific situations where phone use poses genuine risks to safety or fairness. Education about responsible phone use is probably more effective than blanket legislation in most cases.}
\clearpage
\begin{center}{\color{topicblue}\sffamily\bfseries Topic 2: Smiling}\end{center}
\vspace{10pt}
\forecastcuecard{Describe an occasion when many people were smiling}{%
  \item \textbf{When it happened}
  \item \textbf{Who you were with}
  \item \textbf{What happened}
  \item And explain why most people were smiling
}

\forecastsample[1]{If you asked me to pinpoint the happiest I’ve seen my university campus, I’d immediately take you back to last October during our Founder’s Day cultural festival. Usually, these events are a bit formal and stiff, but this particular night turned into something completely unexpected.\par The atmosphere was already buzzing because we were all crammed into the main auditorium for the talent show finale. We were expecting the usual lineup: some acoustic guitar covers, a K-pop dance group, maybe a heartfelt ballad. But then, this senior guy named Minh walked onto the stage. Now, you have to understand, Minh is known for being incredibly quiet—the kind of guy who sits at the back of the lecture hall and barely speaks. So when he grabbed the mic, the room went dead silent, mostly out of confusion.\par Then, he started roasting the entire faculty. And I don’t mean mean-spirited insults; I mean he launched into this brilliant, spot-on impersonation of our strictest professors. He mimicked the way our Philosophy teacher paces back and forth when he forgets his point, and the specific high-pitched "Ahem!" our Dean makes when he’s annoyed.\par It broke the house down. I looked around, and it wasn't just polite chuckles; people were practically falling out of their chairs. It was that belly-laugh kind of joy that brings tears to your eyes. I think what made everyone smile so hard was just the sheer relief of it. We spend so much time stressed about grades and deadlines, terrified of these authority figures, and suddenly, here was one of us turning that fear into comedy. It felt like a massive collective exhale. For those twenty minutes, nobody was worrying about tomorrow’s assignments; we were just a community, united by the sheer absurdity of student life.}
\clearpage
\begin{forecastvocabtable}
\forecastvocabrow{Pinpoint}{v}{/ˈpɪnpɔɪnt/}{to identify something exactly}{chỉ ra chính xác}
\forecastvocabrow{Stiff}{adj}{/stɪf/}{formal and not relaxed}{cứng nhắc}
\forecastvocabrow{Buzzing}{adj}{/ˈbʌzɪŋ/}{full of excitement or activity}{náo nhiệt}
\forecastvocabrow{Roasting}{n}{/ˈrəʊstɪŋ/}{(slang) making fun of someone}{châm chọc (hài hước)}
\forecastvocabrow{Spot-on}{adj}{/ˌspɒt ˈɒn/}{completely accurate}{chính xác tuyệt đối}
\forecastvocabrow{Broke the house down}{idiom}{/brəʊk ðə haʊs daʊn/}{made the audience laugh extremely hard}{làm khán phòng vỡ òa}
\forecastvocabrow{Collective exhale}{n phrase}{/kəˈlektɪv eksˈheɪl/}{everyone relaxing at the same time}{sự thở phào nhẹ nhõm chung}
\end{forecastvocabtable}

\clearpage
\begin{center}{\color{topicblue}\sffamily\bfseries Smiling}\end{center}
\vspace{10pt}
\forecastqa{Do people smile more when they are younger or older?}{From my observations, I'd say younger people, particularly children, tend to smile more freely and frequently. Kids smile at the smallest things – a funny face, a new toy, or just playing with friends. As people age, life's responsibilities and challenges seem to make them more reserved with their expressions. However, I've noticed that elderly people often rediscover that joy and smile quite readily, perhaps because they've learned not to sweat the small stuff. Middle-aged adults, who are often juggling careers and family pressures, probably smile the least. So it's not exactly linear – it's more like a U-shaped curve across a lifetime.}
\forecastqa{Do you think people who like to smile are more friendly?}{Generally speaking, yes, though it's not always a perfect correlation. People who smile often do tend to come across as more approachable and warm, which are hallmarks of friendliness. In Vietnamese culture especially, smiling is seen as a sign of politeness and openness. However, I've met people who smile constantly but are actually quite superficial or even manipulative – it's like a mask they wear. Conversely, some genuinely kind people just have more serious resting faces. So while there's often a connection between smiling and friendliness, you can't judge someone's character solely based on how much they smile. Actions speak louder than facial expressions.}
\forecastqa{Why do most people smile in photographs?}{It's become a social convention, hasn't it? We're conditioned from childhood to "say cheese" and smile for the camera. I think people smile in photos because they want to project happiness and create positive memories – even if they're not feeling particularly joyful in that moment, they want to look back on happy times. There's also social pressure; a serious or neutral expression in photos might be interpreted as being unhappy or unfriendly. In the age of social media, this has intensified because people curate their online image to show their best, happiest selves. Personally, I sometimes find it exhausting to force a smile for every photo, but I do it anyway because that's what's expected.}
\forecastqa{Do women smile more than men? Why?}{In my experience, yes, women do tend to smile more frequently than men, though this is largely due to social conditioning rather than any inherent difference. From a young age, girls are often encouraged to be pleasant, agreeable, and emotionally expressive, while boys are taught to be more stoic and reserved. Women might also smile more as a social strategy – to appear friendly, to defuse tension, or even as a defensive mechanism in uncomfortable situations. I've read that women sometimes smile when they're nervous or trying to be polite, even when they don't feel like it. Men, on the other hand, face less social pressure to constantly appear cheerful and can get away with more neutral expressions without being judged as unfriendly.}
\clearpage
\begin{center}{\color{topicblue}\sffamily\bfseries Topic 3: Pride and Rewards}\end{center}
\vspace{10pt}
\forecastcuecard{Describe a time when you felt proud of a family member}{%
  \item \textbf{When it happened}
  \item \textbf{Who the person is}
  \item \textbf{What the person did}
  \item And explain why you felt proud of him/her
}

\forecastsample[1]{I don't say this enough, but I am incredibly proud of my little sister. The specific moment that solidified this feeling happened about six months ago, but to understand why it mattered, you need to know that she used to be painfully shy—like, "hiding behind my mom’s leg" shy.\par She’s in her final year of high school now, and somehow, she decided to sign up for the National English Debate Championship. When she told me, I honestly thought she was joking. But she was dead serious. For months, I watched her transform her bedroom into a war room, covered in sticky notes and research papers about artificial intelligence ethics.\par The pride really hit me during the finals in Ho Chi Minh City. My parents and I were in the audience, holding our breath as she walked up to the podium. She looked so small against the massive backdrop of the stage. But when she opened her mouth, that shy little girl was gone. She spoke with this fierce, articulate authority that I had never seen before. She didn't just read a script; she actually listened to her opponents, dismantled their arguments piece by piece, and rebuilt her own case on the fly.\par Watching her up there, controlling the room and commanding respect from judges who were twice her age, I actually got choked up. It wasn't just about her winning—which she did, by the way—but about witnessing the sheer grit it took for her to conquer her biggest fear. It made me realize that while I was busy trudging through my university assignments, my little sister had quietly grown into this powerhouse of a young woman. It was a humbling, beautiful moment that completely shifted the dynamic of our relationship.}
\clearpage
\begin{forecastvocabtable}
\forecastvocabrow{Solidified}{v}{/səˈlɪdɪfaɪd/}{made something strong or certain}{củng cố}
\forecastvocabrow{Painfully shy}{adj phrase}{/ˈpeɪnfəli ʃaɪ/}{extremely shy}{ngại ngùng đến đau lòng}
\forecastvocabrow{Dead serious}{adj phrase}{/ded ˈsɪəriəs/}{completely serious}{cực kỳ nghiêm túc}
\forecastvocabrow{War room}{n}{/ˈwɔː ruːm/}{(metaphor) a room used for intense planning}{phòng chiến lược}
\forecastvocabrow{On the fly}{idiom}{/ɒn ðə flaɪ/}{while something is happening; without preparation}{ứng biến ngay lập tức}
\forecastvocabrow{Choked up}{adj}{/tʃəʊkt ʌp/}{feeling like you might cry}{nghẹn ngào}
\forecastvocabrow{Powerhouse}{n}{/ˈpaʊəhaʊs/}{a person with great energy or ability}{người tràn đầy năng lượng/tài năng}
\end{forecastvocabtable}

\clearpage
\begin{center}{\color{topicblue}\sffamily\bfseries Pride and Rewards}\end{center}
\vspace{10pt}
\forecastqa{When would parents feel proud of their children?}{Parents typically feel proud when their children achieve significant milestones or demonstrate good character. Academic achievements are a big one in Vietnamese culture – getting into a good university, winning competitions, or excelling in exams. But beyond academics, parents feel proud when their children show kindness, responsibility, or overcome challenges. My parents, for instance, expressed pride not just when I got into university, but also when I started volunteering and showing initiative in helping others. I think parents also feel proud during those small everyday moments when they see their values reflected in their children's behavior – like when a child shows respect to elders or helps a sibling without being asked.}
\forecastqa{Should parents reward children? Why and how?}{Yes, I believe rewards can be beneficial when used appropriately. They provide positive reinforcement and help children understand that effort and good behavior have consequences. However, the type of reward matters enormously. Instead of always using material rewards like toys or money, which can create a transactional mindset, parents should also use praise, quality time, or special privileges. For example, my parents would sometimes reward good grades by taking me to my favorite restaurant or letting me choose a family activity, which created positive associations without being purely materialistic. The key is ensuring children understand they're being rewarded for genuine effort and improvement, not just outcomes, so they develop intrinsic motivation rather than just chasing external rewards.}
\forecastqa{Is it good to reward children too often? Why?}{Definitely not. Over-rewarding can be counterproductive and create several problems. Firstly, children might start expecting rewards for everything, even basic responsibilities or behaviors that should be intrinsic. This undermines the development of internal motivation – they'll only do things if there's something in it for them. Secondly, constant rewards can diminish their value; if you get a prize for everything, nothing feels special anymore. I've seen this with some younger cousins who've been over-rewarded – they've become quite entitled and struggle with tasks that don't come with immediate gratification. Rewards should be reserved for genuine achievements or exceptional effort, not for meeting basic expectations. Otherwise, you're setting children up for disappointment when they enter the real world where not everything comes with a reward.}
\forecastqa{On what occasions would adults be proud of themselves?}{Adults tend to feel proud when they accomplish goals they've worked hard towards, whether that's professional achievements like promotions or personal milestones like buying a house or completing a marathon. Overcoming adversity is another big one – successfully navigating a difficult period, breaking a bad habit, or recovering from failure. As a university student on the cusp of adulthood, I feel proud when I manage my responsibilities well, like balancing studies with part-time work or helping my family financially. I think adults also feel proud of more subtle achievements that others might not notice – like maintaining their composure in a stressful situation, making an ethical decision despite pressure, or simply being there for someone who needed support. It's often about living up to your own standards and values.}
\clearpage
\begin{center}{\color{topicblue}\sffamily\bfseries Topic 4: Possessions and Toys}\end{center}
\vspace{10pt}
\forecastcuecard{Describe something that you can't live without (not a computer/phone)}{%
  \item \textbf{What it is}
  \item \textbf{What you do with it}
  \item \textbf{How it helps you in your life}
  \item And explain why you can't live without it
}

\forecastsample[1]{If you took away my phone and laptop, I’d survive, but if you took away my physical planner, I would probably disintegrate into chaos within twenty-four hours. It’s a simple, black leather notebook that I carry everywhere, and honestly, it’s the only thing keeping my life from falling apart.\par In a world that’s obsessed with digital apps and cloud syncing, sticking to pen and paper might seem archaic, but for me, it’s psychological anchoring. There is something about the physical act of writing down a deadline that makes it real in a way that typing it into Google Calendar never does. When I scribble "Microeconomics Essay due Friday" in red ink, my brain actually registers the urgency.\par But it’s grown into more than just a to-do list. Over the last few years of university, it’s become this chaotic scrapbook of my life. If you flip through the pages, you won’t just see homework; you’ll see doodles from boring lectures, ticket stubs from movies I saw with friends, and random ideas for essays that hit me at 2 AM. It’s where I break down my overwhelming workload into manageable little boxes that I can physically check off. That little dopamine hit I get from striking a line through a finished task? That is genuinely what gets me through exam season.\par So, it’s not really about the object itself. It’s about the sense of control it gives me. University life feels like being in a constant storm of information and responsibilities, and this notebook is basically my anchor. Without it, I'd just be drifting, missing appointments and forgetting deadlines. It is my external hard drive, my therapist, and my personal assistant all rolled into one.}
\clearpage
\begin{forecastvocabtable}
\forecastvocabrow{Disintegrate}{v}{/dɪsˈɪntɪɡreɪt/}{break into small parts/lose control}{tan rã}
\forecastvocabrow{Archaic}{adj}{/ɑːˈkeɪɪk/}{very old or old-fashioned}{cổ lỗ sĩ}
\forecastvocabrow{Psychological anchoring}{n phrase}{/ˌsaɪkəˈlɒdʒɪkl ˈæŋkərɪŋ/}{making you feel mentally stable}{neo giữ tâm lý}
\forecastvocabrow{Registers the urgency}{v phrase}{/ˈredʒɪstəz ði ˈɜːdʒənsi/}{notices that something is important/urgent}{nhận thức được sự cấp bách}
\forecastvocabrow{Scrapbook}{n}{/ˈskræpbʊk/}{a book with stuck-in photos/cuttings}{cuốn sổ lưu niệm}
\forecastvocabrow{Dopamine hit}{n phrase}{/ˈdəʊpəmiːn hɪt/}{a feeling of pleasure}{cảm giác sung sướng (do hormone dopamine)}
\forecastvocabrow{Drifting}{adj}{/ˈdrɪftɪŋ/}{moving slowly without direction}{trôi dạt vô định}
\end{forecastvocabtable}

\clearpage
\begin{center}{\color{topicblue}\sffamily\bfseries Possessions and Toys}\end{center}
\vspace{10pt}
\forecastqa{Why do all children like toys?}{Toys serve multiple developmental purposes, which is why they're universally appealing to children. Fundamentally, toys provide entertainment and fun, which is crucial for childhood happiness. But beyond that, they're tools for learning and imagination – through play, children explore concepts, develop motor skills, and create narratives. Toys also offer a sense of ownership and control in a world where children have limited autonomy. A child's toy is something that's entirely theirs to command and manipulate, which is empowering. Additionally, I think toys provide comfort and companionship, especially stuffed animals or dolls that become trusted friends. From an evolutionary perspective, play with toys mimics real-world activities and helps children practice skills they'll need as adults.}
\forecastqa{Do you think it is good for a child to always take his or her favourite toy with them all the time?}{It depends on the child's age and the context. For very young children, having a comfort object like a favorite toy can provide security and help them cope with new or stressful situations – it's a transitional object that eases separation anxiety. However, if an older child is overly dependent on a toy and can't function without it, that might indicate an issue that needs addressing. There's also the practical concern of losing or damaging the toy, which could be traumatic. I think it's healthy to gradually encourage independence from such objects as children mature. Perhaps allowing the toy at bedtime or during particularly challenging situations, but not constantly. It's about finding a balance between providing comfort and fostering resilience.}
\forecastqa{Why are children attracted to new things (such as electronics)?}{Children are naturally curious and drawn to novelty – it's how they learn about the world. New things, especially electronics, offer stimulation and instant gratification that traditional toys might not provide. Modern electronics like tablets or gaming devices are designed to be engaging with bright colors, sounds, and interactive elements that capture attention. There's also a social aspect; children want what their peers have, and electronics have become status symbols even among young kids. Additionally, I think children pick up on adult behavior – they see us constantly on our phones and laptops, so naturally, they want to imitate that. Electronics also offer a sense of sophistication and "growing up," which appeals to children's desire to be more mature.}
\forecastqa{Why do some grown-ups hate to throw out old things (such as clothes)?}{There are several psychological and practical reasons. Sentimentality is a big factor – old possessions often carry memories and emotional significance. That worn-out jacket might remind someone of their university days or a special trip. There's also the "just in case" mentality; people hold onto things thinking they might need them someday, even if they haven't used them in years. For some, it's about frugality or growing up with scarcity – they can't bear to waste something that still has potential use. I've noticed my parents are like this, having grown up in less affluent times. There's also decision fatigue; it's easier to keep something than to make the decision to discard it. And honestly, some people just form attachments to their possessions and struggle with letting go.}
\forecastqa{Is the way people buy things affected? How?}{Absolutely, consumer behavior has transformed dramatically, especially with technology and social media. Online shopping has made purchasing incredibly convenient – you can buy almost anything with a few clicks and have it delivered to your door. This ease has probably increased impulse buying. Social media influencers and targeted advertising have also changed how people discover and decide on purchases; we're constantly exposed to products through our feeds. There's also more emphasis on reviews and peer recommendations – before buying anything significant, I always check online reviews and ask friends for opinions. The rise of e-wallets and digital payments has made transactions feel less "real," which might lead to overspending. Additionally, there's growing awareness about sustainability and ethical consumption, so some people are more mindful about what they buy and from whom.}
\clearpage
\begin{center}{\color{topicblue}\sffamily\bfseries Topic 5: Relaxation and Exercise}\end{center}
\vspace{10pt}
\forecastcuecard{Describe your favorite place in your house where you can relax}{%
  \item \textbf{Where it is}
  \item \textbf{What it is like}
  \item \textbf{What you enjoy doing there}
  \item And explain why you feel relaxed at this place
}

\forecastsample[1]{You’d probably expect me to say my bedroom, but honestly, my room is where all my stress lives—it’s where my desk and my textbooks are. So, my actual sanctuary is the tiny balcony right outside my glass sliding door.\par I live in a student apartment, so when I say "balcony," I’m being generous. It’s basically a concrete slab barely big enough for two people to stand on. But I’ve spent the last two years turning it into my own little escape pod. I scavenged a rusty camping chair from a secondhand shop, covered it in cushions, and filled every inch of floor space with potted plants. I’ve probably killed more succulents than I’ve kept alive, but the ones that survived make it feel like a mini jungle.\par This is where I go when the noise of the world gets too loud. Hanoi can be incredibly chaotic—the traffic, the construction, the constant hum of the city. But at night, when I step out there and close the door behind me, the noise fades into this white noise background that’s actually quite soothing. I string up these warm yellow fairy lights, make a cup of tea, and just sit. No phone, no assignments, no "what am I doing with my life" panic spirals.\par It’s the only place where I allow myself to be completely unproductive. Inside, I always feel guilty if I’m not studying. But out there? That’s neutral territory. Whether I’m watching the neighbors cook dinner across the alleyway or just staring at the smoggy night sky, it’s the one spot where I can just breathe. That twenty minutes of solitude is usually the only thing recharging my battery enough to go back inside and face the textbooks again.}
\clearpage
\begin{forecastvocabtable}
\forecastvocabrow{Generous}{adj}{/ˈdʒenərəs/}{kind; (here) describing something as better than it really is}{hào phóng (ở đây nghĩa là nói quá lên)}
\forecastvocabrow{Escape pod}{n}{/ɪˈskeɪp pɒd/}{distinct safe place}{khoang thoát hiểm (nghĩa bóng: nơi trốn tránh)}
\forecastvocabrow{Scavenged}{v}{/ˈskævɪndʒd/}{collected from discarded items}{lượm lặt, tìm kiếm}
\forecastvocabrow{Mini jungle}{n phrase}{/ˈmɪni ˈdʒʌŋɡl/}{a small place with many plants}{khu rừng nhỏ}
\forecastvocabrow{Panic spirals}{n phrase}{/ˈpænɪk ˈspaɪrəlz/}{getting more and more anxious}{vòng xoáy hoảng loạn}
\forecastvocabrow{Neutral territory}{n phrase}{/ˈnjuːtrəl ˈterətri/}{a place belonging to no side}{vùng trung lập}
\forecastvocabrow{Solitude}{n}{/ˈsɒlɪtjuːd/}{the state of being alone}{sự cô đơn (tích cực)}
\end{forecastvocabtable}

\clearpage
\begin{center}{\color{topicblue}\sffamily\bfseries Relaxation and Exercise}\end{center}
\vspace{10pt}
\forecastqa{Why is it difficult for some people to relax?}{I think modern life has created a culture where being busy is glorified, and relaxation is sometimes seen as laziness or wasted time. Many people, especially students and working professionals, feel guilty when they're not being productive. There's also the issue of constant connectivity – with smartphones, we're always reachable and always aware of what we should be doing, which makes it hard to truly switch off. Anxiety and stress can create a vicious cycle where you need to relax but your worried mind won't let you. Some people also simply haven't learned how to relax; they don't have hobbies or practices that help them unwind. For me personally, even when I have free time, I sometimes find myself restlessly thinking about assignments or future responsibilities rather than enjoying the present moment.}
\forecastqa{What are the benefits of doing exercise?}{The benefits are extensive, both physical and mental. Obviously, exercise improves physical health – it strengthens your cardiovascular system, builds muscle, helps maintain a healthy weight, and boosts your immune system. But the mental benefits are equally important, especially for students like me dealing with academic pressure. Exercise releases endorphins, which improve mood and reduce stress and anxiety. I always feel clearer-headed and more focused after a workout. It also improves sleep quality, which is crucial for cognitive function and overall wellbeing. Beyond individual benefits, exercise can be social – playing sports or going to the gym with friends combines fitness with social connection. Long-term, regular exercise reduces the risk of chronic diseases and contributes to longevity and quality of life.}
\forecastqa{Do people in your country exercise after work?}{It's becoming increasingly common, especially in urban areas like Hanoi and Ho Chi Minh City. You'll see people jogging around lakes, practicing tai chi in parks, or going to gyms in the evening. However, I'd say it's still not as widespread as in some Western countries. Many Vietnamese people work long hours and have family responsibilities, so exercise often takes a backseat. There's also a cultural aspect – traditionally, physical labor was seen as work rather than recreation, so the concept of exercising for health is relatively modern. Younger generations are definitely more fitness-conscious, influenced by social media and global health trends. But overall, I'd estimate that maybe 30-40\% of working adults exercise regularly, with higher rates among younger, urban professionals.}
\forecastqa{What is the place where people spend most of their time at home?}{In Vietnamese homes, it's typically the living room or, increasingly, bedrooms. The living room is traditionally the heart of the home where families gather to watch TV, eat meals, and socialize. However, with younger people spending more time on personal devices and studying, bedrooms have become multifunctional spaces – for sleeping, studying, working, and entertainment. In my dorm, I definitely spend most of my time in my bedroom because it's my only private space. For families, the dining area is also significant, especially during meal times which are important for Vietnamese family culture. Kitchens are used frequently but usually just for cooking rather than lingering. I think the pandemic also shifted patterns, with more people working from home and spending time in home offices or study areas.}
\clearpage
\begin{center}{\color{topicblue}\sffamily\bfseries Topic 6: Movies}\end{center}
\vspace{10pt}
\forecastcuecard{Describe a movie you watched recently that you felt disappointed about}{%
  \item \textbf{When it was}
  \item \textbf{Why you didn't like it}
  \item \textbf{Why you decided to watch it}
  \item And explain why you felt disappointed about it
}

\forecastsample[1]{I’m usually pretty easy to please when it comes to movies—give me some popcorn and a decent plot, and I’m happy. But last month, I dragged a group of friends to see this new Vietnamese rom-com called "Yêu Nhầm Bạn Thân," and I practically spent the whole walk home apologizing to them.\par I fell for the marketing hook, line, and sinker. The trailer was everywhere on TikTok, promising this witty, modern take on the "best friends falling in love" trope. It starred these two actors who are currently the "it couple" of Vietnamese cinema, so I assumed the chemistry would be electric. I was expecting something charming, maybe a bit cheesy but heartwarming—perfect for a post-exam brain detox.\par Instead, what we got was two hours of second-hand embarrassment. The script felt like it was written by an AI that had watched too many bad American sitcoms. The dialogue was so unnatural; nobody in real life talks the way they did. They were trying so hard to be "quirky" and "gen Z" that it just came off as painful.\par But the worst part was the complete lack of logic. Characters would make these baffling decisions just to force a misunderstanding, simply because the plot needed conflict. There was this one scene where they broke up over a lost cat, which made absolutely no sense. I sat there in the dark theater, checking my watch, literally cringing at the screen. It felt like such a wasted opportunity because the production value was high—the cinematography was gorgeous, the outfits were trendy—but it was all style and zero substance. Walking out, I didn't feel uplifted; I just felt annoyed that I’d wasted my limited student budget on something that felt so hollow.}
\clearpage
\begin{forecastvocabtable}
\forecastvocabrow{Hook, line, and sinker}{idiom}{/hʊk laɪn ənd ˈsɪŋkə(r)/}{completely (believed a trick/lie)}{tin sái cổ (cắn câu)}
\forecastvocabrow{The "it couple"}{n phrase}{/ði ɪt ˈkʌpl/}{the most popular couple at the moment}{cặp đôi đình đám}
\forecastvocabrow{Brain detox}{n phrase}{/breɪn ˈdiːtɒks/}{doing something simple to rest the brain}{"thải độc" cho não}
\forecastvocabrow{Second-hand embarrassment}{n phrase}{/ˌsekənd hænd ɪmˈbærəsmənt/}{feeling embarrassed for someone else}{ngại dùm người khác}
\forecastvocabrow{Came off as}{phrasal verb}{/keɪm ɒf əz/}{appeared to be}{có vẻ như}
\forecastvocabrow{Baffling}{adj}{/ˈbæflɪŋ/}{impossible to understand}{khó hiểu, kỳ quặc}
\forecastvocabrow{Style and zero substance}{phrase}{/staɪl ənd ˈzɪərəʊ ˈsʌbstəns/}{looks good but has no meaning}{tốt nước sơn mà không tốt gỗ}
\end{forecastvocabtable}

\clearpage
\begin{center}{\color{topicblue}\sffamily\bfseries Movies}\end{center}
\vspace{10pt}
\forecastqa{Do you believe movie reviews?}{I take them with a grain of salt. Professional critic reviews can provide valuable insights about cinematography, storytelling, and thematic elements that I might miss. However, critics sometimes have different priorities than regular viewers – they might appreciate artistic merit while I just want to be entertained. I find user reviews on platforms like IMDb or Rotten Tomatoes more relatable because they reflect ordinary viewers' experiences. That said, I'm aware that reviews can be subjective and influenced by personal taste, hype, or even manipulation. I usually look at the overall consensus rather than individual reviews, and I try to read both positive and negative ones to get a balanced perspective. Ultimately, reviews inform my decision but don't dictate it – I've enjoyed movies that were panned by critics and disliked some that were highly praised.}
\forecastqa{What are the different types of films in your country?}{Vietnamese cinema has diversified quite a bit in recent years. Romantic comedies are probably the most popular and commercially successful – light-hearted films about love and relationships that appeal to young audiences. We also have historical and war films, often depicting Vietnam's struggles for independence, which carry cultural and educational significance. Horror films have gained traction, particularly ones incorporating Vietnamese folklore and supernatural beliefs. There's a growing indie film scene producing more artistic, experimental works that tackle social issues. Action films and crime thrillers are emerging too, though they're still developing compared to Hollywood or Korean productions. Family dramas remain popular, especially during holiday seasons. What's interesting is seeing Vietnamese filmmakers blend traditional storytelling with modern techniques and global influences.}
\forecastqa{Are historical films popular in your country? Why?}{They have a dedicated audience but aren't as mainstream as romantic comedies or horror films. Historical films, particularly those about the Vietnam War or the struggle against colonialism, hold cultural importance and are often supported by the government for their educational value. Older generations tend to appreciate them more because they connect to their lived experiences or family histories. For younger people like myself, historical films can feel heavy or preachy, especially if they prioritize propaganda over storytelling. However, when done well – with good production values, compelling narratives, and nuanced characters – they can be quite engaging. I think there's potential for historical films to become more popular if filmmakers find ways to make them relevant and emotionally resonant for contemporary audiences without sacrificing historical accuracy.}
\forecastqa{Do you think films with famous actors or actresses are more likely to become successful films?}{Generally, yes, but it's not a guarantee. Famous actors bring built-in audiences – their fans will watch the film regardless of the plot or quality. They also attract media attention and marketing opportunities, which increases visibility. In Vietnam, having a popular actor can significantly boost box office numbers, especially for romantic comedies or commercial films. However, I've seen films with big-name stars flop because the story was weak or the execution was poor. Conversely, some films with lesser-known actors have become sleeper hits through word-of-mouth and genuine quality. I think star power gets people into theaters, but it's the film's actual merit that determines long-term success and cultural impact. Ultimately, the best scenario is combining talented, popular actors with strong scripts and skilled direction.}
\forecastqa{Why are Japanese animated films so popular?}{Japanese animation, particularly from Studio Ghibli and directors like Makoto Shinkai, has achieved global popularity for several reasons. Firstly, the storytelling is often sophisticated and emotionally resonant, appealing to both children and adults. Films like "Spirited Away" or "Your Name" deal with complex themes – identity, loss, environmentalism – in accessible ways. The animation quality is exceptional, with meticulous attention to detail and beautiful artistry that's visually stunning. Japanese anime also offers diverse genres and styles, from whimsical fantasy to intense drama, so there's something for everyone. Culturally, there's a certain mystique and aesthetic appeal to Japanese animation that feels fresh compared to Western animation. The music and soundtracks are often memorable and enhance the emotional impact. For Vietnamese audiences specifically, anime has been popular since the 90s, so there's nostalgia and familiarity that drives continued interest.}
\end{document}
